\section{Introduction}
\label{sec:introduction}

Influence maximization (IM) is a canonical sequential optimization problem on social and information networks: given a diffusion graph and a seed budget, the goal is to select a small set of nodes whose activation maximizes expected influence spread. The classical greedy view is conceptually simple, but each greedy step can require comparing conditional marginal gains for many candidates, and accurately estimating these gains through Monte Carlo simulation or scalable sampling-based influence oracles remains computationally expensive on large graphs.

Learning offers an appealing shortcut. A graph model can predict which candidates are promising, or directly estimate the marginal gain of adding a candidate under the current seed set. However, a sequential IM solver needs more than an accurate predictor on average. The relevant ranking changes after every selected seed because of overlap and diminishing returns; a candidate that looks strong globally can become redundant under the current state. Consequently, a locally plausible prediction error can remove the true best candidate from consideration and propagate through later decisions. This makes a fully learned replacement for the influence oracle difficult to trust under distribution shift or predictor failure.

We take a different perspective: a learned marginal predictor should be treated as \emph{fallible advice} to a classical optimizer, not as an oracle that must always be correct. The key question is therefore not ``Can a GNN replace influence estimation?'', but rather ``When is learned advice useful enough to reduce oracle work, and how much verification is needed before acting on it?'' We instantiate this perspective with a learning-augmented IM framework that combines full-graph candidate screening, state-aware marginal proposals, an audited residual trust test, progressive candidate and sample verification, and a classical fallback (Figure~\ref{fig:overview}). When predictions are useful, the learned model concentrates expensive evaluation on a small head of plausible candidates. When predictions are unreliable, the audit expands real-oracle computation and can recover classical behavior.

Our current evidence supports three central claims. First, marginal prediction must be explicitly state aware: supervision that exposes the same candidate under different seed states substantially improves conditional ranking sensitivity. Second, prediction alone is insufficient for reliable seed selection, whereas selective verification recovers much of the lost decision quality. Third, an audited progressive solver yields the desired learning-augmented behavior: useful advice reduces expensive verification, while corrupted or random advice increases oracle effort and activates fallback rather than causing catastrophic downstream degradation. Importantly, this behavior persists when the method receives the full influence graph and uses an independent screening stage to construct a compact candidate set.

\paragraph{Contributions.}
Our contributions are threefold:
\begin{itemize}
    \item We formulate learned marginal-gain prediction for sequential IM as a \emph{fallible advice interface}, highlighting why ordinary offline prediction quality is insufficient to guarantee reliable downstream decisions.
    \item We propose a learning-augmented solver that integrates full-graph screening, state-aware marginal proposals, residual auditing, progressive verification, and fallback to classical influence evaluation.
    \item We evaluate the resulting quality--efficiency--robustness tradeoff, including controlled predictor degradation and full-graph inputs, and show that oracle computation adapts to advice quality rather than assuming a uniformly reliable learned model.
\end{itemize}
